\section{Raw experiments, task fibers and paid acquisition}
\label{sec:setup}
All spaces are standard Borel. A positive normalized causal kernel, a preparation and a fixed bounded terminal audit $Z\in H$ are given, where $H\subset\R^k$ is compact and convex. The model $\theta\in\Theta$ is fixed along a trajectory and unavailable to the observer. The hidden state may include its microscopic past. The complete visible filtration contains every reported failure, every previously used private coin, and every recorded physical cost. A fresh routing coin is conditionally independent of the future audit. Every legal terminating procedure has a stopping time $1\le\tau<\infty$ almost surely under every admitted model; the infimum is restricted to such procedures.

At raw call $t$ a physical command is selected from the already available information. The nonnegative reported cost $c_t$ satisfies $c_t\ge c_*>0$. A predictable flag $A_t\in\{0,1\}$ says whether this call is productive. Only a productive report permits commitment. Its visible task mark is $X_t\in\mathsf X$; an additional report $V_t$ can be informative and may be used by arbitrary competing procedures. The audit is executed after commitment, with no undeclared intervening transition. Preparation and audit costs not included in $c_t$ are displayed separately when a total budget is used.

Fix a task baseline $B_0\ge0$. After the entire current report let
\[
 \mu=\E[Z\mid\mathcal H],\qquad s=\E[\|Z-\mu\|^2\mid\mathcal H]-B_0.
\]
We assume $s\ge0$ on stop-allowed histories. For each mark $x$, let $\mathcal F_x\subset H\times[0,S]$ be the closure of the attainable pairs $(\mu,s)$ over all admitted models and legal full histories with that mark, excluding zero-evidence histories. It is a task fiber, not a set chosen by the encoder. Conditional versions on their actual exposure laws suffice. For the pointwise fiber criterion below we require nonempty compact fibers and Borel functions
\[
 \underline s(x)=\inf_{(u,v)\in\mathcal F_x}v,
 \qquad g_x(r)=\sup_{(u,v)\in\mathcal F_x}\{v+\|u-r\|^2\}
\]
for each $r$ in a fixed countable dense subset of $H$. In applications an explicitly verified compact invariant family can replace the actual fibers for a \emph{sufficient} certificate; failure on a larger family is not a necessity statement about the actual experiment.

The unique minimizer and its value are
\begin{equation}\label{eq:canonical-mark}
 f_*(x)=\argmin_{r\in H}g_x(r),\qquad w_*(x)=g_x(f_*(x)).
\end{equation}
They are Borel, as proved in Theorem~\ref{thm:fiber}. Its raw calibration criterion is
\begin{equation}\label{eq:fiber-ratio}
 w_*(x)\le K\underline s(x),\qquad 1\le K<\infty.
\end{equation}
It implies the uniform comparison
\begin{equation}\label{eq:task-compare}
 k_-\{w(x)+\|f(x)-r\|^2\}
 \le \E[\|Z-r\|^2\mid\mathcal H]-B_0
 \le k_+\{w(x)+\|f(x)-r\|^2\},
\end{equation}
with $(f,w)=(f_*,w_*)$, $k_-=(3K)^{-1}$ and $k_+=2$. Equivalently one may verify any bounded Borel $f:\mathsf X\to H$, $w\ge0$ and positive $k_-,k_+$ satisfying \eqref{eq:task-compare}. For example
\begin{equation}\label{eq:ordinary-calibration}
 a_vw\le s\le A_vw,\qquad \|\mu-f\|^2\le a_bw
\end{equation}
gives $k_-=[\max\{2,(1+2a_b)/a_v\}]^{-1}$ and $k_+=\max\{2,A_v+2a_b\}$. All comparisons are after $V_t$; hence they survive selection using $V_t$ and an independent fresh stopping coin.

\subsection*{The acquired law and its conditional envelopes}
Fix a known reference law $\nu$ of the mark produced by a declared, actually executable reference preparation and acquisition. Its acquisition is not repeated for free during the online task. A known reference member of a model family may be used; it is not the unknown true member. Assume constants $0<\ell\le1\le u<\infty$ such that, for every model and every surviving productive history and command,
\begin{equation}\label{eq:mass-upper}
 \law_\theta(X_t\mid\mathcal H_{t-1},\text{command})\le u\nu.
\end{equation}
This is a measure inequality on all Borel sets. The actual conditional laws can differ for every history.

A parameter-independent acquisition program $\mathcal G$ has $B$ live labels, a designated return label, and compatible physical commands. From that return label it runs until a productive report; after rejection it returns to that label without resetting the hidden physical state. It stores no other history, seed, real posterior or clock. At its productive calls require also
\begin{equation}\label{eq:mass-lower}
 \law_\theta(X_t\mid\mathcal H_{t-1},\text{command})\ge\ell\nu.
\end{equation}
The current mark and the chosen terminal cell can be evaluated from its live label and the current report. Every attempt completes almost surely and has conditional expected raw cost at most $\kappa_+$, uniformly over its starting history and model. These are verified properties of an actual finite-state program, not properties assigned to an abstract state graph.

\subsection*{A lower effort certificate}
There is a bounded Borel potential $V(h)$ on raw continued histories, independent of the true model, equal to zero initially and immediately after every productive call. For every legal pre-call history and command,
\begin{equation}\label{eq:effort-drift}
 \E_\theta[c_t+V(\mathcal H_t)-V(\mathcal H_{t-1})
            \mid\mathcal H_{t-1},\text{command}]
 \ge\kappa_- A_t,\qquad \kappa_->0.
\end{equation}
$V$ may be negative. At commitment its value is zero. A limiting version with uniformly integrable stopped potentials is possible, but boundedness is the hypothesis used here. This certificate constrains \emph{all} admissible strategies, not just $\mathcal G$.

\begin{lemma}[Verification of an acquisition program]\label{lem:navigator}
The upper attempt-cost property follows if a nonnegative function $W(s,h)$ on the product of its live labels and raw predictive histories satisfies $W(s_0,h)\le\kappa_+$ at each return history and
\[
 \E_\theta[c_t+W(S_t,\mathcal H_t)\one_{\{\text{attempt continues}\}}
      \mid S_{t-1},\mathcal H_{t-1}]
 \le W(S_{t-1},\mathcal H_{t-1})
\]
for every surviving attempt history. The proof potential is not stored by the machine.
\end{lemma}
\begin{proof}
Stop at the earlier of attempt completion and $n$ further calls. Conditional telescoping, with the nonnegative remaining potential retained, bounds expected cost by the initial $W$. Monotone convergence bounds the full sum. Since each call costs at least $c_*$, the expected number of calls is finite and completion is almost sure. This argument is uniform over the starting history. It assumes neither regeneration nor a positive completion hazard on every raw call.
\end{proof}

Let $R_\cp^\Theta(N,m)$ be the infimum of worst-model risk over arbitrary legal history-dependent procedures with a fixed dictionary of at most $m$ terminal values and $\sup_\theta\E_\theta\sum_{t\le\tau}c_t\le N$. Let $R_\aut^\Theta(N,M)$ impose a total alphabet of at most $M$ disjoint live and terminal labels. Actions depend only on the live label; routing uses that label and the current command and report. The routing class at this level consists of Borel stochastic kernels. This is an alphabet-constrained model, not a claim that every Borel rule has finite program description or finite arithmetic cost. Those resources require a computable cell certificate as in Proposition~\ref{prop:finite-cells}, or a direct finite-bit construction as in Section~\ref{sec:numerics}. Infeasible risks are $+\infty$. The same $B_0$ and the same audit occur in both risks.
